\setcounter{paragraphcounter}{1}

\section{Arbeitsdienstordnung}

\subsection{Präambel}
Diese Arbeitsdienstordnung schafft klare, faire und transparente Regelungen, um eine gerechte Verteilung der Aufgaben sicherzustellen und den reibungslosen Ablauf aller Vereinsaktivitäten zu gewährleisten.

\subsection{§ \theparagraphcounter\ Erlass, Änderung, Aufhebung und Bekanntmachung}
\begin{enumerate}
    \item Diese Arbeitsdienstordnung kann jederzeit durch den Vereinsausschuss erarbeitet und beschlossen werden.

    \item Änderungen werden durch den Beschluss des Vereinsausschusses wirksam und allen Mitgliedern bekannt gemacht.
\end{enumerate}

\refstepcounter{paragraphcounter}
\subsection{§ \theparagraphcounter\ Zweck der Arbeitsdienste}
\begin{enumerate}
    \item Die Arbeitsdienste dienen der Durchführung von Vereinsveranstaltungen, sowie der Unterstützung des Vereinsbetriebs.

    \item Jedes Mitglied ist im Rahmen seiner Mitgliedschaft verpflichtet, die zugeteilten Arbeitsdienste ordnungsgemäß zu leisten.

     \item Zur Ableistung von Arbeitsdiensten verpflichtet sind ausschließlich ordentliche Mitglieder vom vollendeten 15. Lebensjahr bis zum vollendeten 65. Lebensjahr. Mitglieder außerhalb dieses Altersbereiches, außerordentliche und Ehrenmitglieder sind von der Arbeitsdienstpflicht befreit.
\end{enumerate}

\refstepcounter{paragraphcounter}
\subsection{§ \theparagraphcounter\ Zuständigkeit}
\begin{enumerate}
    \item Für die Organisation, Planung und Durchführung der Arbeitsdienste ist der Festwart verantwortlich.

    \item Der Festwart kann zur Erfüllung seiner Aufgaben durch Mitglieder des Ausschusses unterstützt werden.

    \item Der Vereinsausschuss kann ergänzende Regelungen oder Weisungen erlassen, sofern dies zur ordnungsgemäßen Durchführung der Arbeitsdienste erforderlich ist.
\end{enumerate}

\refstepcounter{paragraphcounter}
\subsection{§ \theparagraphcounter\ Dienstplan}
\begin{enumerate}
    \item Der Festwart erstellt für alle Vereinsveranstaltungen, insbesondere die jährlichen Feste, einen Dienstplan.

    \item Die Einteilung der Mitglieder erfolgt in einem rollierenden System jeweils für den Auf- bzw. Abbau und für die Durchführung des Festes.

    \item Der Festwart achtet auf eine faire und gleichmäßige Verteilung der Arbeitsdienste.

    \item Der fertige Dienstplan ist den Mitgliedern spätestens zwei Wochen vor Festbeginn in geeigneter Form mitzuteilen.
\end{enumerate}

\refstepcounter{paragraphcounter}
\subsection{§ \theparagraphcounter\ Erscheinen und Abmeldung}
\begin{enumerate}
    \item Die Mitglieder sind verpflichtet, pünktlich zu den eingeteilten Arbeitsdiensten zu erscheinen.

    \item Kann ein Mitglied den zugewiesenen Dienst aus zwingenden Gründen nicht antreten, so ist dies dem Festwart unverzüglich mitzuteilen.

    \item Das verhinderte Mitglied hat eigenständig für geeigneten Ersatz zu sorgen und dem Festwart diesen mitzuteilen.

    \item Der Festwart ist dafür verantwortlich zu prüfen, ob alle eingeteilten Mitglieder anwesend sind und gegebenenfalls Ersatzkräfte einzuteilen.
\end{enumerate}

\refstepcounter{paragraphcounter}
\subsection{§ \theparagraphcounter\ Dokumentation der Arbeitsstunden}
\begin{enumerate}
    \item Die geleisteten Arbeitsstunden werden vom Festwart oder einem Vorstandsmitglied dokumentiert.

    \item Die Dokumentation muss vollständig, nachvollziehbar und zeitnah erfolgen.

    \item Die erfassten Stunden dienen als Grundlage für die Anrechnungen auf das Sprunggeld gemäß § 7 Abs. 6 der Vereinssatzung.
\end{enumerate}

\refstepcounter{paragraphcounter}
\subsection{§ \theparagraphcounter\ Versäumnisse}
\begin{enumerate}
    \item Nicht geleistete oder unentschuldigte Arbeitsdienste gelten als Versäumnis.

    \item Wiederholte Verstöße gegen die Arbeitsdienstordnung können vereinsrechtliche Maßnahmen nach sich ziehen.
\end{enumerate}

\refstepcounter{paragraphcounter}
\subsection{§ \theparagraphcounter\ Inkrafttreten}
\begin{enumerate}
    \item Diese Arbeitsdienstordnung tritt nach Beschluss der Ausschusssitzung vom dd. MMMM YYYY in Kraft. \\
    Sie ersetzt alle vorherigen Versionen
\end{enumerate}